\chapter{Einleitung}
Dieses Experiment untersucht die Lebensdauer des angeresgten $\frac{5}{2}^+$-Niveaus von $^{133}_{55}\mathrm{Cs}$. 
Dieses Wird durch die Zerfalls Kette von $^{133}_{56}\mathrm{Ba}$ bestimmt, bei der ein angeregter Zustand des Cäsiumkerns entsteht und anschließend durch Emission eines $\gamma$-Quants in einen Energetisch tieferen Zustand übergeht.
Die dabei emittierte Strahlung wird mit zwei Szintillationsdetektoren registriert. Um die für die Messung relevanten Kernübergänge eindeutig identifizieren zu können, erfolgt zunächst eine Energiekalibration der Detektoren sowie die Bestimmung ihrer Energieauflösung. 
Auf dieser Grundlage werden geeignete Energiefenster für die nachfolgende Koinzidenzmessung festgelegt.

Zur Untersuchung des zeitlichen Verhaltens des Messaufbaus wird $^{22}_{11}\mathrm{Na}$ entstehenden Koinzidenzphotonen bestimmt. 
Danach werden die Zeitdifferenzen zwischen den aufeinanderfolgenden $\gamma$-Übergängen des $^{133}_{56}\mathrm{Ba}$-Zerfalls gemessen.
Aus der resultierenden Zeitverteilung lässt sich die Lebensdauer des $\frac{5}{2}^{+}$-Zustands von $^{133}_{55}\mathrm{Cs}$ ermitteln.